\documentclass[11pt]{article}
\usepackage{amsmath,amssymb,amsthm}
\usepackage[margin=1in]{geometry}
\usepackage{booktabs}
\usepackage{algorithm}
\usepackage{algpseudocode}
\usepackage[colorlinks=true,linkcolor=blue,citecolor=blue,urlcolor=blue]{hyperref}

\newtheorem{theorem}{Theorem}[section]
\newtheorem{proposition}[theorem]{Proposition}
\newtheorem{lemma}[theorem]{Lemma}
\newtheorem{corollary}[theorem]{Corollary}
\theoremstyle{definition}
\newtheorem{definition}[theorem]{Definition}
\newtheorem{remark}[theorem]{Remark}
\newtheorem{question}[theorem]{Open question}

\setlength{\emergencystretch}{2em}

\newcommand{\B}{\mathbf{B}}
\newcommand{\bb}{\mathbf{b}}
\newcommand{\dd}{\mathbf{d}}
\newcommand{\A}{\mathbf{A}}
\newcommand{\vv}{\mathbf{v}}
\newcommand{\x}{\mathbf{x}}
\newcommand{\kk}{\mathbf{k}}
\newcommand{\gv}{\mathbf{g}}
\newcommand{\R}{\mathbb{R}}
\newcommand{\T}{\mathbb{T}}
\newcommand{\Z}{\mathbb{Z}}
\newcommand{\C}{\mathbb{C}}
\newcommand{\Sph}{\mathbb{S}}
\newcommand{\Tr}{\mathsf{T}}
\newcommand{\Trr}{\mathsf{T}_3}
\newcommand{\Pl}{\mathsf{P}}
\newcommand{\Wm}{\mathsf{W}}
\newcommand{\Zf}{\mathsf{Z}}
\newcommand{\Aop}{\mathsf{A}}
\newcommand{\Minv}{\mathsf{M}^{-1}}
\newcommand{\dive}{\operatorname{div}}
\newcommand{\curl}{\operatorname{curl}}
\newcommand{\supp}{\operatorname{supp}}
\newcommand{\tr}{\operatorname{tr}}
\newcommand{\rank}{\operatorname{rank}}

\title{A minimum-norm Gauss--Newton method for carrier-free\\
continuation of constant-$|\B|$ fields}
\author{A.~Mallet\\
\small Space Sciences Laboratory, University of California, Berkeley}
\date{August 2026}

\begin{document}
\maketitle

\begin{abstract}
We document the algorithm behind the v2 (``carrier-free'') solver stack of the
\texttt{jax\_constantB} package: a minimum-norm Gauss--Newton continuation for
divergence-free unit vector fields on the periodic box, in which the
divergence multiplier is eliminated analytically. The elimination reduces each
sweep to a single scalar normal equation, which we show is exactly the
minimum-norm correction expressed in vector-potential variables; it preserves
$\nabla\cdot\B$ to round-off by construction, and its Galerkin (strict
$2/3$-rule) discretisation has an \emph{exact} Jacobian, so that Newton's
method converges quadratically with the explicitly computable constant
$\tfrac12$. We derive the symbol of the resulting normal operator,
$(1+k^2)^{-s}\sin^2\theta(\hat\kk,\B)$, which identifies the removable part
of the conditioning (a diagonal preconditioner) and the intrinsic part
(field-aligned near-null directions). We then give the constraints used to
pin genuine three-dimensionality, an experimental descent that seeks the
smoothest state at fixed pinned energies, and the continuation driver.
Measured convergence, iteration counts and timings are collected in
Table~\ref{tab:measured}.
\end{abstract}

\section{The problem}\label{sec:problem}

\subsection{Physical motivation}

A magnetic field of constant magnitude and vanishing divergence, together with
the velocity perturbation $\delta\mathbf u=\pm\,\delta\B/\sqrt{4\pi\rho}$, is
an exact nonlinear solution of ideal incompressible MHD: a
\emph{spherically polarised} (arc-polarised, large-amplitude) Alfv\'en state.
The solar wind is observed to be full of such states, including the
large-amplitude ``switchbacks'' in which the field direction reverses locally
relative to the mean. Whether genuinely three-dimensional states of this kind
can be continued to switchback amplitude while remaining smooth --- or whether
they are forced to develop discontinuities --- is the scientific question this
solver exists to answer. Everything below is purely kinematic: we are
constructing divergence-free maps $\B:\T^3\to\Sph^2$, and the MHD content
enters only through that identification.

\subsection{Formulation}

Work on the flat torus $\T^3=[0,2\pi)^3$, so wavenumbers are integer vectors
$\kk\in\Z^3$. We seek $\B:\T^3\to\R^3$ with
\begin{equation}\label{eq:constraints}
\nabla\cdot\B=0,\qquad |\B(\x)|=1 \quad\text{for all }\x .
\end{equation}
Write the constraint map
\begin{equation}\label{eq:F}
F(\B)=\bigl(\nabla\cdot\B,\; q\bigr),\qquad
q:=\tfrac12\bigl(|\B|^2-1\bigr),
\end{equation}
so that \eqref{eq:constraints} is $F(\B)=0$. The system is
\emph{underdetermined}: two scalar equations for three unknown components at
every point. Solutions therefore come in large families, and a solver must
choose which member it moves toward; that choice --- here, a weighted
minimum-norm choice --- is a modelling decision, not an implementation
detail, and it is the mechanism by which paths of different smoothness reach
equal amplitude.

The states of interest are continued from a uniform mean field
$\bar\B$ (with $|\bar\B|<1$ along the way; see
Corollary~\ref{cor:mean}) plus an arbitrary divergence-free seed
$\bb=\nabla\times\mathbf a$, injected in small increments. This is the
\emph{carrier-free} formulation: no one-dimensional arc-polarised carrier wave
is used as a backbone, in contrast with the v1 stack. Decomposing
$\B=\bar\B+\bb$ with $\langle\bb\rangle=0$, the mean of the quadratic
constraint gives the exact identity
\begin{equation}\label{eq:energy}
\langle|\bb|^2\rangle = 1-|\bar\B|^2 ,
\end{equation}
$\langle\cdot\rangle$ denoting the box average: the fluctuation energy is
determined by the mean field alone.

\subsection{Notation}\label{sec:notation}

Throughout, all fields are real; $\hat f(\kk)$ denotes Fourier coefficients
with the convention $\langle f,g\rangle=(2\pi)^{-3}\!\int f\cdot g
=\sum_\kk \hat f(\kk)^*\!\cdot\hat g(\kk)$, and $k^2=|\kk|^2$. The operators
used below are all Fourier multipliers and are collected here:

\begin{center}
\small
\begin{tabular}{@{}lp{0.80\textwidth}@{}}
\toprule
$\Tr$ & strict $2/3$ band projector on scalars, $\widehat{\Tr f}(\kk)=
\mathbf 1[\,|k_i|<N_i/3\ \forall i\,]\,\hat f(\kk)$\\
$\Trr$ & the same projector applied componentwise to a stacked $3$-vector
field\\
$\Pl$ & Leray (divergence-free) projector,
$\widehat{\Pl\vv}(\kk)=(\mathsf I-\hat\kk\hat\kk^{\!\top})\hat\vv(\kk)$ with
$\hat\kk=\kk/|\kk|$; $\Pl=\mathsf I$ at $\kk=0$\\
$\Wm^{-2}$ & Sobolev weight, $\widehat{\Wm^{-2}f}=(1+k^2)^{-s}\hat f$, with
$s=\texttt{smooth}\ge0$; $\Wm^{2}$ its inverse\\
$\Zf$ & freeze mask: spectral indicator, $0$ on pinned bins (and their
Hermitian partners), $1$ elsewhere\\
$\Aop$ & the scalar normal operator of \eqref{eq:normal-mu}\\
$\Minv$ & the CG preconditioner, $\Minv=\Tr(1+k^2)^{s}$\\
\bottomrule
\end{tabular}
\end{center}

\noindent
$\Tr,\Trr,\Pl,\Wm^{\pm2},\Zf$ are all diagonal in $\kk$ (scalar multipliers
except $\Pl$, which is a matrix multiplier), hence mutually commuting,
self-adjoint, and $\Tr^2=\Tr$, $\Pl^2=\Pl$, $\Zf^2=\Zf$. We use these facts
repeatedly without comment. $N_i$ is the grid size on axis $i$; grids are
written $N_x^2\times N_z$ when $N_x=N_y$.

\section{The minimum-norm Gauss--Newton step}\label{sec:mngn}

Linearising \eqref{eq:F} about $\B$ gives the Jacobian
\begin{equation}\label{eq:J}
J\dd=\bigl(\nabla\cdot\dd,\;\B\cdot\dd\bigr),
\end{equation}
whose $L^2$ adjoint follows from
$\langle J\dd,(\lambda,\mu)\rangle
=\langle \lambda,\nabla\cdot\dd\rangle+\langle\mu,\B\cdot\dd\rangle
=\langle \dd, -\nabla\lambda+\mu\B\rangle$ (integrating by parts on the
torus, no boundary term):
\begin{equation}\label{eq:JT}
J^{\!\top}(\lambda,\mu)=-\nabla\lambda+\mu\B .
\end{equation}
Because $J$ has more columns than rows the Newton correction is not unique. We
select it by a weighted minimum-norm condition: with $s\ge0$ fixed, take
\begin{equation}\label{eq:minnorm}
\dd=\arg\min\ \tfrac12\|\Wm\dd\|^2
\quad\text{subject to}\quad J\dd=-F(\B),
\qquad
\|\Wm\dd\|^2=\sum_\kk(1+k^2)^{s}\,|\hat\dd(\kk)|^2 .
\end{equation}
The Lagrangian $\tfrac12\|\Wm\dd\|^2-\langle(\lambda,\mu),J\dd+F\rangle$
stationarises at $\Wm^2\dd=J^{\!\top}(\lambda,\mu)$, so
\begin{equation}\label{eq:step-lm}
\dd=\Wm^{-2}J^{\!\top}(\lambda,\mu),\qquad
\bigl(J\,\Wm^{-2}J^{\!\top}\bigr)(\lambda,\mu)=-F(\B).
\end{equation}
For $s=0$ this is the classical $L^2$ minimum-norm (Gauss--Newton for
underdetermined systems) step; for $s>0$ it is the \emph{smoothest} correction
that cancels the linearised residual, and increasing $s$ biases the
continuation path toward smoother members of the solution family. The v1
solver solves \eqref{eq:step-lm} for the pair $(\lambda,\mu)$ by matrix-free
conjugate gradients. Section~\ref{sec:reduction} eliminates $\lambda$.

The outer iteration is
\begin{equation}\label{eq:sweep}
\B\ \longleftarrow\ \B+\dd ,
\end{equation}
which we call a \emph{sweep}. Since $F$ is quadratic, the exact second-order
remainder is available in closed form and will be used in
Section~\ref{sec:galerkin}:
\begin{equation}\label{eq:remainder}
q(\B+\dd)=q(\B)+\B\cdot\dd+\tfrac12|\dd|^2,\qquad
\nabla\cdot(\B+\dd)=\nabla\cdot\B+\nabla\cdot\dd .
\end{equation}

\section{Elimination of \texorpdfstring{$\lambda$}{lambda}: the
vector-potential form}\label{sec:reduction}

This section is the centrepiece. The divergence row of \eqref{eq:step-lm} can
be solved for $\lambda$ in closed form, independently of the weight, leaving a
single scalar equation for $\mu$; the resulting step is precisely the
minimum-norm correction written in vector-potential variables.

\subsection{The Schur reduction}

\begin{proposition}[$\lambda$-elimination]\label{prop:schur}
Let $\B$ be divergence-free, so that the first component of $F(\B)$ vanishes,
and let $\Wm^{-2}$ be any Fourier multiplier of the form $w(k^2)$. Then the
solution of \eqref{eq:step-lm} satisfies
\begin{equation}\label{eq:lambda}
\hat\lambda(\kk)=-\,\frac{i\,\kk\cdot\widehat{(\mu\B)}(\kk)}{k^2}
\qquad(\kk\neq0),
\end{equation}
independently of $w$, and the correction reduces to
\begin{equation}\label{eq:step-mu}
\boxed{\ \dd=\Wm^{-2}\Pl(\mu\B)\ }
\end{equation}
with $\Pl$ the Leray projector. The remaining (Schur-complement) normal
equation is the scalar problem
\begin{equation}\label{eq:normal-mu}
\Aop\mu:=\B\cdot\Wm^{-2}\Pl(\mu\B)=-q .
\end{equation}
\end{proposition}

\begin{proof}
The divergence row of \eqref{eq:step-lm} reads
$\nabla\cdot\bigl[\Wm^{-2}(-\nabla\lambda+\mu\B)\bigr]=-\nabla\cdot\B=0$. In
Fourier variables, with $\widehat{\nabla\lambda}=i\kk\hat\lambda$ and
$\widehat{\nabla\cdot\vv}=i\kk\cdot\hat\vv$, this is
\[
w(k^2)\Bigl[k^2\hat\lambda(\kk)
+i\,\kk\cdot\widehat{(\mu\B)}(\kk)\Bigr]=0 ,
\]
which for $\kk\ne0$ gives \eqref{eq:lambda} whenever $w(k^2)\neq0$; the
common factor $w$ cancels, which is why the elimination is weight-independent.
Substituting into $\dd=\Wm^{-2}(-\nabla\lambda+\mu\B)$,
\[
\hat\dd(\kk)=w(k^2)\Bigl[\widehat{(\mu\B)}(\kk)
-\frac{\kk\,\bigl(\kk\cdot\widehat{(\mu\B)}(\kk)\bigr)}{k^2}\Bigr]
=w(k^2)\,\bigl(\mathsf I-\hat\kk\hat\kk^{\!\top}\bigr)\widehat{(\mu\B)}(\kk).
\]
At $\kk=0$ the divergence row is vacuous, $\lambda$ is undetermined (a
constant, which contributes nothing), and $\hat\dd(0)=\widehat{(\mu\B)}(0)$,
consistent with the convention $\Pl=\mathsf I$ at $\kk=0$. This is
\eqref{eq:step-mu}. The second row of \eqref{eq:step-lm},
$\B\cdot\dd=-q$, is then \eqref{eq:normal-mu}.
\end{proof}

\begin{proposition}[symmetry and positivity]\label{prop:spd}
$\Aop$ is self-adjoint and positive semi-definite; indeed
$\Aop=G^{\!\top}G$ with $G\mu:=\Wm^{-1}\Pl(\mu\B)$, and
$\langle\mu,\Aop\mu\rangle=\|\Wm^{-1}\Pl(\mu\B)\|^2$. Its null space is
$\{\mu:\Pl(\mu\B)=0\}$; in particular $\Aop\,1=|\B|^2=1$ at a solution, so
constants are not null.
\end{proposition}

\begin{proof}
$\langle\mu',\Aop\mu\rangle=\langle\mu'\B,\Wm^{-2}\Pl(\mu\B)\rangle$. Since
$\Wm^{-2}$ and $\Pl$ commute, are self-adjoint, and $\Pl^2=\Pl$, we may write
$\Wm^{-2}\Pl=(\Wm^{-1}\Pl)^{\!\top}(\Wm^{-1}\Pl)$, giving
$\langle\mu',\Aop\mu\rangle=\langle G\mu',G\mu\rangle$. For the last
statement, $\Pl\B=\B$ because $\B$ is divergence-free.
\end{proof}

Proposition~\ref{prop:spd} is what licenses conjugate gradients on
\eqref{eq:normal-mu}, and it survives every modification made below (band
projection, freeze masks), because each is a commuting self-adjoint
projection.

\subsection{The vector-potential reading}

The correction \eqref{eq:step-mu} is divergence-free by construction, hence a
curl. The following identifies the potential and shows that the reduction is
not merely algebraic bookkeeping: it is the minimum-norm step taken in the
variables $\A$, with $\B=\bar\B+\nabla\times\A$.

\begin{proposition}[potential form]\label{prop:potential}
For any field $\vv$ with $\langle\vv\rangle=0$, let
$\A_\vv:=\nabla\times(-\nabla^2)^{-1}\vv$, i.e.\
$\hat\A_\vv(\kk)=i\kk\times\hat\vv(\kk)/k^2$. Then $\A_\vv$ is in the Coulomb
gauge, $\nabla\cdot\A_\vv=0$, and
\begin{equation}\label{eq:curlcurl}
\nabla\times\A_\vv=\Pl\vv .
\end{equation}
Consequently $\dd=\nabla\times\dd\A$ with
$\dd\A=\Wm^{-2}\nabla\times(-\nabla^2)^{-1}(\mu\B)$: the step
\eqref{eq:step-mu} is the curl of the (weighted) minimum-norm Coulomb-gauge
potential generated by $\mu\B$.
\end{proposition}

\begin{proof}
Coulomb gauge is immediate from
$\kk\cdot(\kk\times\hat\vv)=0$. For \eqref{eq:curlcurl} use
$\nabla\times\nabla\times\vv=-\nabla^2\vv+\nabla(\nabla\cdot\vv)$, applied to
$(-\nabla^2)^{-1}\vv$:
\[
\nabla\times\A_\vv
=\nabla\times\nabla\times(-\nabla^2)^{-1}\vv
=\vv-\nabla(\nabla\cdot(-\nabla^2)^{-1}\vv),
\]
whose Fourier transform is
$\hat\vv-\kk(\kk\cdot\hat\vv)/k^2=\Pl\hat\vv$.
\end{proof}

\begin{remark}[which norm is minimised]\label{rem:norms}
For a mean-free divergence-free field $\dd$ with Coulomb-gauge potential
$\dd\A$ we have $\hat\dd=i\kk\times\hat{\dd\A}$ with
$\hat{\dd\A}\perp\kk$, hence the exact modewise relation
\begin{equation}\label{eq:normshift}
|\hat{\dd\A}(\kk)|=|\hat\dd(\kk)|/|\kk| .
\end{equation}
Minimising $\|\dd\A\|_{L^2}$ is therefore minimising
$\sum_\kk|\hat\dd|^2/k^2$, a homogeneous $H^{-1}$ norm on $\dd$: working in
potential variables shifts the smoothing index by one. The weighted families
correspond as $s_\A\leftrightarrow s_\B+1$, exactly for the homogeneous weight
$k^{2s}$ and up to the low-$\kk$ difference between $(1+k^2)$ and $k^2$ for
the inhomogeneous weight used in the code. We keep the $\B$-side convention
$(1+k^2)^{-s}$ throughout, so that $s=0$ means the plain $L^2$ minimum-norm
step and no operator has a singularity at $\kk=0$.
\end{remark}

The practical significance is that the unknown of a sweep is a single scalar
field $\mu$, that the divergence constraint is discharged analytically rather
than iteratively, and that the CG operator applies $\Pl$ --- one batched
transform pair --- instead of a gradient and a divergence.

\subsection{Two corollaries}

\begin{corollary}[divergence is preserved, not repaired]\label{cor:div}
Every step \eqref{eq:step-mu} satisfies $\nabla\cdot\dd=0$ exactly (to
round-off), so $\nabla\cdot\B$ is invariant along the whole solve. Measured:
$\max|\nabla\cdot\B|<10^{-13}$ per sweep at $64^2\times128$. Conversely, a
pre-existing divergence error is \emph{not} removed by any number of sweeps.
\end{corollary}

\begin{proof}
$\kk\cdot\widehat{\Pl\vv}(\kk)=0$ by construction, and $\Wm^{-2}$ is a scalar
multiplier. Invariance in both directions follows from
\eqref{eq:remainder}.
\end{proof}

This is a deliberate trade. In the coupled $(\lambda,\mu)$ system the
divergence row carries the right-hand side $-\nabla\cdot\B$ and the solver
spends iterations removing a divergence error that the linear projector
removes exactly and for free. The v2 convention is therefore: \emph{project
once at entry},
\begin{equation}\label{eq:project}
\B\ \longleftarrow\ \Trr\,\Pl\,\B \qquad\text{(mean preserved)},
\end{equation}
and never again. The API separates the two (\texttt{project} versus
\texttt{gn}) so that the assumption of Proposition~\ref{prop:schur} is visible
at the call site rather than buried.

\begin{corollary}[the mean field]\label{cor:mean}
The $\kk=0$ component of a step is
$\dd\bar\B=\langle\mu\B\rangle$. Zeroing that bin (equivalently, including
$\kk=0$ in the freeze mask $\Zf$) holds $\bar\B$ fixed for the entire
continuation. Fixed-mean growth is only possible for $|\bar\B|<1$: if
$|\bar\B|=1$ then $\B\equiv\bar\B$.
\end{corollary}

\begin{proof}
The first claim is $\Pl=\mathsf I$ at $\kk=0$. For the second, Cauchy--Schwarz
(Jensen) gives $1=|\langle\B\rangle|^2\le\langle|\B|^2\rangle=1$ with equality
if and only if $\B$ is constant almost everywhere; this is
\eqref{eq:energy} with $\langle|\bb|^2\rangle=0$.
\end{proof}

Corollary~\ref{cor:mean} is enforced at the driver level: \texttt{--fix-mean}
is refused with an explicit message when $|\bar\B|\ge1-10^{-9}$, because the
requested continuation then provably does not exist.

\section{Galerkin discretisation}\label{sec:galerkin}

\subsection{\texorpdfstring{The strict $2/3$ band}{The strict 2/3 band}}

Fix a grid $N_x\times N_y\times N_z$ and the \emph{retained band}
\begin{equation}\label{eq:band}
\mathcal K_N:=\bigl\{\kk\in\Z^3:\ |k_i|<N_i/3\ \text{for }i=1,2,3\bigr\},
\end{equation}
with band projector $\Tr$ (componentwise $\Trr$). The strict inequality is
load-bearing. With the inclusive cutoff $|k_i|\le N_i/3$, the product of two
modes sitting exactly on the cutoff has a component at $2N_i/3$, which aliases
to $2N_i/3-N_i=-N_i/3$: exactly onto the retained band edge. Orszag's
exactness argument then fails for the quadratic nonlinearity, and the
Galerkin claims of this section become false. (In the code this is the
$-10^{-12}$ guard in the axis mask; the comparison must not be left to
floating-point luck when $N_i$ is a multiple of $3$.)

The unknown is a band-limited field, $\B=\Trr\B$, and the discretised
constraint map is
\begin{equation}\label{eq:FG}
F_G(\B)=\bigl(\nabla\cdot\B,\;\Tr q\bigr),\qquad q=\tfrac12(|\B|^2-1),
\end{equation}
the divergence of a band-limited field being band-limited already. The
reported residual is $\max|\Tr q|$.

\begin{proposition}[Orszag exactness]\label{prop:orszag}
Let $u,v$ be band-limited, $\supp\hat u,\supp\hat v\subset\mathcal K_N$, and
let $\widehat{uv}^{\,\mathrm{disc}}$ denote the DFT of their pointwise product
sampled on the $N$-grid. Then
$\widehat{uv}^{\,\mathrm{disc}}(\kk)=\widehat{uv}(\kk)$ for all
$\kk\in\mathcal K_N$: the retained-band equations of \eqref{eq:FG} are exact
continuum equations, carrying no discretisation error whatsoever.
\end{proposition}

\begin{proof}
Sampling folds frequencies modulo $N_i$ on each axis:
$\widehat{uv}^{\,\mathrm{disc}}(\kk)=\sum_{\mathbf j\in\Z^3}
\widehat{uv}(\kk+N\mathbf j)$, where $N\mathbf j$ means
$(N_1j_1,N_2j_2,N_3j_3)$. The true product has
$\supp\widehat{uv}\subset\mathcal K_N+\mathcal K_N\subset
\{|k_i|<2N_i/3\}$. For $\kk\in\mathcal K_N$ and $j_i\neq0$,
$|k_i+N_ij_i|\ge N_i-N_i/3=2N_i/3$, so every alias term vanishes.
\end{proof}

\subsection{Tail honesty}

The out-of-band content of $q$ on the \emph{same} grid is the diagnostic
\texttt{tail\_norm}: the rms and max of $q-\Tr q$. Its meaning is the
following.

\begin{proposition}[alias fold, one-sided]\label{prop:tail}
With $\B$ band-limited, write $c_\mathbf m:=\hat q(\mathbf m)$ for the true
(continuum) coefficients of $q$, supported in $\{|m_i|<2N_i/3\}$, and
$d_\kk$ for the discrete coefficients at grid bins $\kk\notin\mathcal K_N$.
Then $d_\kk=\sum_{\mathbf m\equiv\kk\,(\mathrm{mod}\,N)}c_\mathbf m$, each
class containing at most $2$ frequencies per axis and hence at most $8$ in
total, and
\begin{equation}\label{eq:tailbound}
\Bigl(\sum_{\mathbf m\notin\mathcal K_N}|c_\mathbf m|^2\Bigr)^{1/2}
\ \ge\ \frac{1}{\sqrt8}\,
\Bigl(\sum_{\kk\notin\mathcal K_N}|d_\kk|^2\Bigr)^{1/2}.
\end{equation}
\end{proposition}

\begin{proof}
The fold formula is Proposition~\ref{prop:orszag}'s identity at
$\kk\notin\mathcal K_N$; the in-band bins receive no aliases, so all true
tail power lands out of band. On axis $i$, the frequencies of a residue class
lie in an interval of length $4N_i/3$ and are spaced $N_i$ apart, so at most
two are occupied. Cauchy--Schwarz gives $|d_\kk|^2\le8\sum|c_\mathbf m|^2$
over the class, and summing over the (disjoint) classes gives
\eqref{eq:tailbound}.
\end{proof}

Thus \texttt{tail\_norm} measures the genuine unresolved burden of the
continuum problem on a single grid, up to a factor $\sqrt8$ in the safe
direction: it never overstates the true tail by more than that. It can in
principle understate it, if two occupied fold partners cancel; we have not
observed this, and an exact check is cheap when wanted, because
$\supp\hat q\subset\{|k_i|<2N_i/3\}$ lies strictly inside the Nyquist band of
the $2\times$ zero-padded grid, so the padded $q$-spectrum is alias-free
(\texttt{spectra.q\_spectrum}). The point of the Galerkin mode is that the
honest diagnostic is available at working resolution: a collocation residual,
by contrast, is blind to aliasing by construction.

\subsection{The exact Galerkin Jacobian and quadratic convergence}

Differentiating \eqref{eq:FG} in a band-limited direction $\dd$,
\begin{equation}\label{eq:JG}
J_G\dd=\bigl(\nabla\cdot\dd,\ \Tr(\B\cdot\dd)\bigr),
\qquad
J_G^{\!\top}(\lambda,\mu)=\Trr\bigl(-\nabla\lambda+\mu\B\bigr),
\end{equation}
the adjoint being computed in the space of band-limited fields:
$\langle\Tr(\B\cdot\dd),\mu\rangle=\langle\dd,\Trr(\mu\B)\rangle$ for
$\mu=\Tr\mu$. Repeating Proposition~\ref{prop:schur} verbatim (all operators
involved commute) gives the discrete step and normal equation actually
implemented:
\begin{equation}\label{eq:step-G}
\dd=\Wm^{-2}\Zf\,\Pl\,\Trr(\mu\B),
\qquad
\Aop\mu=\Tr\bigl[\B\cdot\Wm^{-2}\Zf\,\Pl\,\Trr(\mu\B)\bigr]=-\Tr q ,
\end{equation}
with $\Zf\equiv1$ when nothing is frozen. Proposition~\ref{prop:spd} applies
unchanged, since $\Tr$ and $\Zf$ are commuting self-adjoint projections:
$\Aop=G^{\!\top}G$ with $G\mu=\Wm^{-1}\Zf\Pl\Trr(\mu\B)$.

\begin{theorem}[exact quadratic sweep]\label{thm:quad}
Let $\B=\Trr\B$ be divergence-free and band-limited, let $\dd$ solve
\eqref{eq:step-G} exactly with $\Zf\equiv1$, and set $\B'=\B+\dd$. Then
$\B'$ is band-limited and divergence-free, and
\begin{equation}\label{eq:quad}
\Tr q(\B')=\tfrac12\,\Tr\bigl(|\dd|^2\bigr),
\qquad\text{hence}\qquad
\bigl\|\Tr q(\B')\bigr\|_\infty\le\tfrac12\|\dd\|_\infty^2 .
\end{equation}
\end{theorem}

\begin{proof}
$\dd$ is band-limited because $\Trr$ appears inside \eqref{eq:step-G}, so
$\Trr\B'=\B'$ and the sweep's truncation is a no-op; it is divergence-free by
Corollary~\ref{cor:div}. Apply $\Tr$ to the exact expansion
\eqref{eq:remainder}: $\Tr q(\B')=\Tr q+\Tr(\B\cdot\dd)+\tfrac12\Tr|\dd|^2$,
and $\Tr(\B\cdot\dd)=-\Tr q$ is exactly the normal equation. Note that
$\Tr(|\dd|^2)$ is itself exact by Proposition~\ref{prop:orszag}.
\end{proof}

Theorem~\ref{thm:quad} is not an asymptotic statement: the new residual
\emph{equals} half the band projection of $|\dd|^2$, with no neglected terms.
Since $\Aop$ has an $O(1)$ symbol (Section~\ref{sec:symbol}), $\|\dd\|$ is
comparable to $\|\Tr q\|$ and the observed contraction constant should be
close to $\tfrac12$. It is: the measured sequences in
Table~\ref{tab:measured}(a) give ratios
$1.3\times10^{-5}/(5.0\times10^{-3})^2=0.52$ and
$8.3\times10^{-11}/(1.3\times10^{-5})^2=0.49$ at $s=0$, and $0.54$, $0.51$ at
$s=1$, before the fourth sweep hits the round-off floor. Four sweeps take the
residual from $5\times10^{-3}$ to machine precision.

\subsection{Relation to the legacy coupled solver}\label{sec:legacy}

The v1 solver is the coupled system \eqref{eq:step-lm} in $(\lambda,\mu)$,
discretised with \emph{collocation} operators as the CG step model while
driving the Galerkin residual $\Tr q$ to zero --- deliberately an inexact
Newton method. That choice was empirically forced: with band-truncated
$(\lambda,\mu)$ operators and a CG iteration budget, the solve exhibited a
limit cycle with a large residual amplification, attributed at the time to
near-null band-edge directions of the projected normal operator. As inexact
Newton it converges only linearly, roughly a factor $1.5$--$2$ per sweep below
the tail scale.

The measurements reported here revise the attribution, and we state the
revision as an empirical finding with its evidence attached. The $\mu$-only
formulation uses the \emph{exact Galerkin} operator \eqref{eq:step-G}, band
projections and all, and no limit cycle occurs at any $s$ or grid we have run;
convergence is quadratic with the constant predicted by
Theorem~\ref{thm:quad}. The distinction is therefore not ``projected versus
unprojected operators'' --- v2 projects --- but the elimination of $\lambda$
together with the removal of a truncated CG budget on a coupled, badly
conditioned system. We attribute the limit cycle to the latter, not to
operator projection per se. We have not built a minimal reproducer isolating
the two effects, so this is an inference from two implementations' behaviour
rather than a proof. The legacy rule of thumb (expect linear convergence, set
\texttt{--res-ok} near $10^{-2}\times$\texttt{gal\_tail\_rms}) is retained for
the legacy solver and does \emph{not} apply to v2, where the residual should
reach the round-off floor in a handful of sweeps.

\section{The normal operator: symbol, conditioning, cost}\label{sec:symbol}

\subsection{Symbol}

\begin{proposition}[frozen-coefficient symbol]\label{prop:sym}
Freeze $\B$ at a constant unit vector $\B_0$ (the leading term of the symbol
for slowly varying $\B$). Then $\Aop$ acts on $\hat\mu(\kk)$ by
multiplication by
\begin{equation}\label{eq:symbol}
a(\kk)=(1+k^2)^{-s}\Bigl(|\B_0|^2-(\B_0\cdot\hat\kk)^2\Bigr)
=(1+k^2)^{-s}\sin^2\theta(\hat\kk,\B_0)\ \in[0,1],
\end{equation}
using $|\B_0|=1$; $\theta$ is the angle between $\kk$ and the local field.
\end{proposition}

\begin{proof}
$\mu\B_0$ has transform $\hat\mu\B_0$; applying
$\Pl$ gives $\hat\mu(\mathsf I-\hat\kk\hat\kk^{\!\top})\B_0$, applying
$\Wm^{-2}$ multiplies by $(1+k^2)^{-s}$, and contracting with $\B_0$ gives
$(1+k^2)^{-s}\hat\mu(|\B_0|^2-(\B_0\cdot\hat\kk)^2)$.
\end{proof}

The symbol factorises into an extrinsic and an intrinsic part.

\paragraph{The $(1+k^2)^{-s}$ factor is removable.} It spans
$1$ to $(1+k_{\max}^2)^{-s}$ across the band. At $64^2\times128$ the band edge
is $|k_i|<21$, so $k^2$ reaches $\sim1.3\times10^3$ and the weighted operator
($s=1$) has a spectral spread of order $10^3$ from this factor alone, on top
of the geometric factor. It is exactly cancelled by the diagonal
preconditioner
\begin{equation}\label{eq:pc}
\Minv=\Tr\,(1+k^2)^{s},
\end{equation}
which costs one scalar transform pair per CG iteration. Table~\ref{tab:measured}(b)
shows the effect: unpreconditioned weighted CG stalls (more than $1500$
iterations, budget-limited at relative residual $\sim10^{-7}$), while
preconditioned CG reaches relative residual $10^{-13}$ in a few hundred.

\paragraph{The $\sin^2\theta$ factor is intrinsic.} It degenerates for
$\kk\parallel\B$: a modulation of $\mu$ along the field produces no admissible
correction at all. The mechanism is transparent in the frozen-coefficient
picture --- if $\B_0=\hat z$ and $\mu=\mu(z)$ then
$\mu\B_0=\nabla\!\int^z\!\mu$ is a pure gradient, annihilated by $\Pl$ --- and
it is the linearised shadow of the rigidity of band-limited constant-$|\B|$
fields (extreme modes forced null and transverse, so that genuinely 3D states
cannot terminate their spectrum): the directions along which the constraint
manifold offers no motion are exactly the field-aligned ones, and they are
never absent. The factor is bounded in $[0,1]$, so it costs conditioning but
not solvability, and it is not preconditionable by any Fourier multiplier.

\subsection{Cost per CG iteration}

Applying $\Aop$ in the form \eqref{eq:step-G} costs: one batched
three-component real transform pair inside $\Pl\Wm^{-2}\Zf\Trr$ (six real
scalar transforms, executed as two batched calls), plus one scalar transform
pair for the band projection of the output ($2$); the preconditioner
\eqref{eq:pc} adds one scalar pair ($2$). The pointwise multiplications
$\mu\B$ and $\B\cdot(\cdot)$ require no transform. The coupled $(\lambda,\mu)$
operator, by comparison, costs $8$ real transforms unweighted and $11$
weighted, plus $2$ for its preconditioner, and carries four fields through the
CG recurrence rather than one. Combined with the iteration counts, this is the
factor $5$--$8$ in wall-clock per sweep recorded in
Table~\ref{tab:measured}(c).

\section{Measurements}\label{sec:measured}

All numbers in Table~\ref{tab:measured} were measured on 2026-08-23 on a CPU
development machine, at $64^2\times128$, starting from an $\epsilon\approx1$
switchback state; the production target is a Kaggle P100, where absolute
timings differ but the ratios do not. CG counts are to relative residual
$10^{-13}$ (the code's stop is $\|r\|^2<10^{-26}\|r_0\|^2$).

\begin{table}[htbp]
\centering
\caption{Measured behaviour of the v2 ($\mu$-only) solver against the legacy
coupled $(\lambda,\mu)$ solver. CPU development machine, $64^2\times128$,
$\epsilon\approx1$ switchback state, 2026-08-23.}
\label{tab:measured}
\medskip
\small
\begin{tabular}{lcccc}
\multicolumn{5}{c}{\textbf{(a) Galerkin Newton residual $\max|\Tr q|$ per sweep}}\\
\toprule
& sweep 1 & sweep 2 & sweep 3 & sweep 4\\
\midrule
$s=0$ & $5.0\times10^{-3}$ & $1.3\times10^{-5}$ & $8.3\times10^{-11}$ & $8.2\times10^{-16}$\\
$s=1$ & $5.1\times10^{-3}$ & $1.4\times10^{-5}$ & $9.9\times10^{-11}$ & $1.0\times10^{-15}$\\
\bottomrule
\end{tabular}

\medskip
\begin{tabular}{lcc}
\multicolumn{3}{c}{\textbf{(b) CG iterations to relative residual $10^{-13}$}}\\
\toprule
system & $s=0$ & $s=1$\\
\midrule
$\mu$-only, preconditioned \eqref{eq:pc} & $244$--$284$ & $291$--$366$\\
$\mu$-only, unpreconditioned & --- & $>1500$ (stall, rel.\ res.\ $\sim10^{-7}$)\\
coupled $(\lambda,\mu)$, Laplacian-block pcg & $916$--$1499$ & not converged in budget\\
\bottomrule
\end{tabular}

\medskip
\begin{tabular}{lccc}
\multicolumn{4}{c}{\textbf{(c) Wall-clock per sweep (CPU) and transform budget}}\\
\toprule
system & $s=0$ & $s=1$ & real transforms per CG iteration\\
\midrule
coupled $(\lambda,\mu)$ + pcg & $23$\,s & $47$\,s & $8$ ($s=0$) / $11$ ($s>0$), $+2$ pcg\\
$\mu$-only + $\Wm^{-2}$ & $4.9$\,s & $5.7$\,s & $6+2$, $+2$ pcg\\
\bottomrule
\end{tabular}

\medskip
\begin{tabular}{ll}
\multicolumn{2}{c}{\textbf{(d) Invariants}}\\
\toprule
$\max|\nabla\cdot\B|$ per sweep & $<10^{-13}$ (Corollary~\ref{cor:div})\\
observed quadratic constant in \eqref{eq:quad} & $0.49$--$0.54$ (predicted $\tfrac12$)\\
\bottomrule
\end{tabular}
\end{table}

\section{Pinning genuine three-dimensionality}\label{sec:pin}

The solution family contains one-dimensional members (arc- and circularly
polarised waves) which are smooth at every amplitude and are strong attractors
for any smoothing-biased path. Since the scientific question concerns
\emph{genuinely three-dimensional} states, the continuation needs a constraint
that provably excludes the degenerate classes.

\begin{definition}\label{def:dim}
Let $S=\supp\hat\B\setminus\{0\}$. The state $\B$ is \emph{one-dimensional
(1D)} if $S$ is contained in a line through the origin, and \emph{planar
(2.5D)} if $S$ is contained in a plane through the origin. Otherwise it is
\emph{genuinely three-dimensional}.
\end{definition}

\begin{lemma}[three non-coplanar pins suffice]\label{lem:pin}
Let $\kk_1,\kk_2,\kk_3\in\Z^3\setminus\{0\}$ with
$\det[\kk_1\,\kk_2\,\kk_3]\neq0$, and suppose the state satisfies
$|\hat\B(\kk_j)|^2=e_j>0$ for $j=1,2,3$. Then $\B$ is neither 1D nor planar.
\end{lemma}

\begin{proof}
The hypothesis puts $\kk_1,\kk_2,\kk_3\in S$. If $S$ lay in a line through the
origin the three would be pairwise parallel, and if $S$ lay in a plane through
the origin the three would be coplanar; either way
$\det[\kk_1\,\kk_2\,\kk_3]=0$, a contradiction.
\end{proof}

A quantitative version uses the $\kk$-space inertia tensor of the
fluctuation, which is also the runtime monitor
(\texttt{spectra.kspace\_inertia}):
\begin{equation}\label{eq:inertia}
\mathsf T:=\sum_{\kk\neq0}|\hat\B(\kk)|^2\,\hat\kk\hat\kk^{\!\top},
\qquad \lambda_1\ge\lambda_2\ge\lambda_3\ge0,
\qquad \tr\mathsf T=1-|\bar\B|^2,
\end{equation}
the trace identity being \eqref{eq:energy}. Then $\rank\mathsf T=1$ iff 1D,
$\rank\mathsf T\le2$ iff planar, and under Lemma~\ref{lem:pin}'s hypothesis
$\mathsf T\succeq\sum_je_j\hat\kk_j\hat\kk_j^{\!\top}\succ0$, because
$\sum_je_j(\hat\kk_j\cdot\vv)^2=0$ forces $\vv$ orthogonal to three
non-coplanar vectors. The driver logs $\lambda_2/\lambda_1$ and
$\lambda_3/\lambda_1$; $(0,0)$ is 1D, $(x,0)$ planar, both positive genuinely
3D.

\begin{remark}[why one scalar pin fails]\label{rem:scalarpin}
A natural cheap alternative is a single scalar overlap constraint
$\langle\B,\bb\rangle=c$ with $\bb$ the 3D seed. It does not pin anything.
Pick any $\kk_0\in\supp\hat\bb$ and consider the family of 1D circularly
polarised states whose support lies on the line $\R\kk_0$; the value of the
single linear functional $\langle\B,\bb\rangle$ varies continuously over that
family (through the relative phase and amplitude of the mode at $\kk_0$) and
sweeps an interval containing $c$ for any $c$ below the Cauchy--Schwarz
ceiling $\|\B\|\,\|\bb\|$. A one-dimensional constraint cannot exclude an
infinite-dimensional degenerate class: the whole prescribed overlap can be
carried by a single mode of a 1D state. Lemma~\ref{lem:pin} avoids this by
constraining three \emph{separate} bin energies, which is a statement about
the support, not about a projection.
\end{remark}

\subsection{Implementation (v1): frozen modes}

The cheap implementation masks the correction off the pinned bins. Let $\Zf$
be the spectral indicator that vanishes on the chosen bins and $1$ elsewhere;
the step is \eqref{eq:step-G}. Two points deserve emphasis.

\emph{First, this is still an exact constrained minimum-norm step.} Solving
\eqref{eq:minnorm} with the additional linear constraints
$\hat\dd(\kk_j)=0$, the stationarity condition becomes
$\Wm^2\dd=J^{\!\top}(\lambda,\mu)+(\mathsf I-\Zf)\xi$ for a multiplier
field $\xi$; imposing $\dd=\Zf\dd$ and using $\Zf(\mathsf I-\Zf)=0$ together
with the commutation of $\Zf$ with $\Wm^{-2}$ and $\Pl$ gives
$\dd=\Zf\Wm^{-2}J^{\!\top}(\lambda,\mu)$, which is \eqref{eq:step-G}. The
$\lambda$-elimination goes through modewise: on retained bins it is
Proposition~\ref{prop:schur}, and on frozen bins the divergence row is
vacuous and $\hat\dd$ vanishes anyway. Symmetry and positivity
(Proposition~\ref{prop:spd}) are unaffected.

\emph{Second, the Hermitian-partner rule.} In the \texttt{rfftn} layout the
$k_z=0$ plane stores both $(k_x,k_y,0)$ and its conjugate partner
$(-k_x,-k_y,0)$ in the same half-array (the $k_z=N_z/2$ plane lies outside the
retained band, so this is the only in-band case). Any per-bin mask must zero
\emph{both} partners; otherwise the inverse real transform silently
symmetrises and the mask leaks a factor-of-two error into the frozen
amplitude.

Under freezing, the pinned amplitudes evolve only through the explicit seed
injection of the continuation driver (Section~\ref{sec:driver}), never through
the solve. The mean-field freeze of Corollary~\ref{cor:mean} is the same
mechanism with $\kk=0$.

\subsection{Implementation (v3): bordered energy rows}\label{sec:erows}

Freezing fixes the pinned \emph{coefficients}; the scientifically meaningful
constraint is on the pinned \emph{energies}, which leaves the phases free.
This is what the code does (\texttt{MuSolver(\dots,\,pins=)} and
\texttt{gn(\dots,\,pin\_targets=)}); the freeze of the previous subsection is
retained only to reproduce the measurement of Table~\ref{tab:pins}.
Define, for a conjugate pair $\{\pm\kk_j\}$ with spectral projector
$\Pl_{\kk_j}$ (a self-adjoint projection on real fields),
\begin{equation}\label{eq:erows}
e_j(\B):=\langle\B,\Pl_{\kk_j}\B\rangle=2|\hat\B(\kk_j)|^2,
\qquad
\gv_j:=\nabla e_j=2\,\Pl_{\kk_j}\B ,
\end{equation}
so that $\delta e_j=\langle\gv_j,\dd\rangle$. Appending these $m$ rows to the
Jacobian,
\begin{equation}\label{eq:Jtilde}
\tilde J\dd=\bigl(\nabla\cdot\dd,\ \Tr(\B\cdot\dd),\
\langle\gv_j,\dd\rangle_{j=1}^m\bigr),
\end{equation}
the minimum-norm step becomes, after the same $\lambda$-elimination,
\begin{equation}\label{eq:bordered-step}
\dd=\Wm^{-2}\Pl\Trr(\mu\B)+\sum_{j=1}^m\nu_j\,\Wm^{-2}\gv_j ,
\end{equation}
with $(\mu,\nu)\in\Tr L^2\times\R^m$ solving the bordered SPD system
\begin{equation}\label{eq:bordered}
\begin{pmatrix}
\Aop & C\\ C^{\!\top} & D
\end{pmatrix}
\begin{pmatrix}\mu\\ \nu\end{pmatrix}
=\begin{pmatrix}-\Tr q\\ \Delta e\end{pmatrix},
\qquad
(C\nu)=\sum_j\nu_j\,\Tr\bigl[\B\cdot\Wm^{-2}\gv_j\bigr],
\quad
D_{jl}=\langle\gv_j,\Wm^{-2}\gv_l\rangle .
\end{equation}
The border is cheap. Each $\gv_j$ is supported on a single conjugate mode
pair, so $\Wm^{-2}\gv_j=(1+|\kk_j|^2)^{-s}\gv_j$ is a scalar multiple (no
transform), $\Pl\gv_j=\gv_j$ (each mode of a divergence-free field is
transverse), and distinct pairs are orthogonal, so
$D=\mathrm{diag}\bigl((1+|\kk_j|^2)^{-s}\|\gv_j\|^2\bigr)$ is diagonal and
positive. One may therefore either run CG on the whole $(\mu,\nu)$ vector, or
eliminate $\nu=D^{-1}(\Delta e-C^{\!\top}\mu)$ and run CG on the rank-$m$
modified operator $\Aop-CD^{-1}C^{\!\top}$, each of whose applications costs
$m$ inner products beyond \eqref{eq:step-G}.

\paragraph{As implemented.} The code takes the first option: one
preconditioned CG on the whole $(\mu,\nu)$ vector, with the $\mu$-block
preconditioner \eqref{eq:pc} unchanged and the $\nu$-block divided by
$D_{jj}$, evaluated as $\langle\gv_j,\Wm^{-2}\Pl\Tr\gv_j\rangle$ so that a
pin whose bin some other mask kills cannot divide by zero. One application of
the bordered operator costs one application of \eqref{eq:step-G} plus $m$
inner products, as advertised: the direction
$\mathsf D=\Wm^{-2}\Pl\Trr(\mu\B+\sum_j\nu_j\gv_j)$ is formed once and then
paired against $\B$ and against each $\gv_j$. Three implementation points
carry the measured behaviour.

\emph{Re-linearisation.} $\gv_j$, $e_j$ and $D_{jj}$ are rebuilt from the
current $\B$ at \emph{every} sweep, not once per solve. The pin residual then
converges quadratically alongside $\Tr q$ (measured: $10^{-13}$ relative
within three or four sweeps, and $\sim10^{-14}$ throughout a continuation)
instead of lagging it.

\emph{Discrete inner product.} Both borders use the same sum over grid points
as the CG, so on an $N_x\!\times\!N_y\!\times\!N_z$ grid the implemented pin
functional is $\mathrm{vol}\cdot e_j$ with $\mathrm{vol}=N_xN_yN_z$; targets
are quoted in those units. This is not cosmetic: a volume mean in the rows
against a sum in the CG destroys the symmetry of \eqref{eq:bordered} and the
pins stop converging. With the convention consistent, the discrete bordered
operator is symmetric to $10^{-14}$ relative, with pins alone, with frozen
bins alone, and with both.

\emph{The $k_z=0$ partner rule} of the previous subsection applies verbatim to
the pin selectors: a row that selects one partner of a conjugate pair and not
the other measures half the energy and steers the wrong quantity.

\paragraph{Measured: energy rows cost what freezing does not.} A frozen bin
cannot participate in cancelling the residual its own presence creates, and
the solver buys feasibility with a spectral cascade. Table~\ref{tab:pins}
compares the two implementations along one continuation (a random seed,
$36^2\times72$, three top modes pinned, the same $\delta\epsilon$ schedule),
against the unpinned control.

\begin{table}[htbp]
\centering
\caption{Frozen coefficients versus bordered energy rows on the same
carrier-free path (random seed, $36^2\times72$, $m=3$ non-coplanar pins,
Galerkin $s=0$). ``tail'' is the Galerkin tail rms, the honest unresolved
burden; $\max|\nabla\B|_F$ is the driver's gradient monitor. The energy rows
land on the unpinned control; the freeze does not.}
\label{tab:pins}
\medskip
\small
\begin{tabular}{lcccc}
\toprule
pinning & $\epsilon$ & $\max|\nabla\B|_F$ & tail & pin rel.\ error\\
\midrule
frozen coefficients \eqref{eq:step-G} & $0.269$ & $2.45$ & $1.2\times10^{-3}$ & ---\\
none (control) & $0.269$ & $0.80$ & $1.7\times10^{-6}$ & ---\\
energy rows \eqref{eq:bordered} & $0.290$ & $0.87$ & $3.8\times10^{-6}$ & $8\times10^{-14}$\\
energy rows \eqref{eq:bordered} & $0.356$ & $1.11$ & $8.3\times10^{-6}$ & $1\times10^{-13}$\\
\bottomrule
\end{tabular}
\end{table}

\begin{remark}[no activation threshold]\label{rem:pinstart}
The freeze needs one: at $\B\approx\bar\B$ the quadratic row reads
$\Tr(\bar\B\cdot\dd)$, so a residual component in a frozen bin can only be
cancelled by $\dd$ in that same bin, and freezing from $\epsilon=0$ makes the
first steps infeasible. An energy row has no such obstruction --- it is a
single scalar equation whose target starts at zero and grows like
$\epsilon^2$, which the seed injection itself satisfies to leading order.
Continuation with the rows active from $\epsilon=0$ is therefore the default,
and the driver offers no delayed activation for them. The natural schedule is
$c_j(\epsilon)=\epsilon^2e_j(\bb)$ with $\bb$ the seed: the energy the
linearised push would put in pin $j$. Because $e_j$ is a grid sum, it must be
recomputed from the seed on each rung of the resolution ladder.
\end{remark}

\section{Smoothest-state descent (experimental)}\label{sec:descent}

The pinned-energy machinery makes a sharper question available: given a
genuinely 3D anchor of prescribed strength, \emph{how much spectral tail is
forced}? Operationally, minimise a smoothness functional over the constraint
manifold at fixed pinned energies.

Take
\begin{equation}\label{eq:R}
R(\B)=\tfrac12\|\Wm\B\|^2=\tfrac12\sum_\kk(1+k^2)^s|\hat\B(\kk)|^2,
\qquad \nabla R(\B)=\Wm^2\B,\quad\text{so}\quad \Wm^{-2}\nabla R(\B)=\B .
\end{equation}

\begin{proposition}[SQP step]\label{prop:sqp}
Minimising $R(\B+\dd)$ subject to the linearised constraints $J\dd=-F(\B)$
gives
\begin{equation}\label{eq:sqp}
\dd=-\B+\Wm^{-2}J^{\!\top}\nu,
\qquad
\bigl(J\Wm^{-2}J^{\!\top}\bigr)\nu=J\B-F(\B) :
\end{equation}
the same normal system as \eqref{eq:step-lm} with a modified right-hand side.
After the $\lambda$-elimination and Galerkin projection this reads
\begin{equation}\label{eq:sqp-mu}
\Aop\mu=\Tr q+1,\qquad \dd=-\B+\Wm^{-2}\Pl\Trr(\mu\B).
\end{equation}
\end{proposition}

\begin{proof}
Stationarity of
$\tfrac12\|\Wm(\B+\dd)\|^2-\langle\nu,J\dd+F\rangle$ gives
$\Wm^2(\B+\dd)=J^{\!\top}\nu$, i.e.\ \eqref{eq:sqp}; substituting into
$J\dd=-F$ gives the stated right-hand side. For \eqref{eq:sqp-mu}: the
divergence row of $J\B-F$ is $\nabla\cdot\B-\nabla\cdot\B=0$, so
Proposition~\ref{prop:schur} applies unchanged, and the quadratic row is
$\Tr(|\B|^2)-\Tr q=\Tr(2q+1)-\Tr q=\Tr q+1$.
\end{proof}

Equivalently, in tangent-gradient form: the $\Wm^{-2}$-metric projection of
$\nabla R$ onto $\ker J$ is $\B-\Wm^{-2}J^{\!\top}m$ with
$(J\Wm^{-2}J^{\!\top})m=J\Wm^{-2}\nabla R=J\B$, so that
$\dd=-(\B-\Wm^{-2}J^{\!\top}m)$ is minus the tangential gradient of $R$. The
two coincide at a converged state, where $F=0$; away from it, \eqref{eq:sqp}
also restores feasibility to first order. In practice we take a damped step
$\B\leftarrow\B+\alpha\dd$ ($\alpha\approx0.3$) followed by a Gauss--Newton
retraction (Section~\ref{sec:galerkin}) back onto the manifold, and accept the
pair only if $R$ decreases while the pinned energies and the residual return
to tolerance.

\begin{remark}[the $s=0$ degeneracy]\label{rem:s0}
At $s=0$ the functional is constant on the manifold, $R=\tfrac12\|\B\|^2$ with
$\|\B\|^2=\langle|\B|^2\rangle=1$, so no descent exists --- and the algebra
reproduces this exactly: with $\Wm^{-2}=\mathsf I$ and $\Tr q=0$,
\eqref{eq:sqp-mu} is $\Aop\mu=1$, which is solved by $\mu\equiv1$ because
$\Aop1=\B\cdot\Pl\B=|\B|^2=1$; then
$\dd=-\B+\Pl\B=-\B+\B=0$. The descent is only meaningful for $s>0$.
\end{remark}

\begin{remark}[pins are mandatory here]\label{rem:pins-descent}
Without the energy rows of \eqref{eq:bordered} the descent has an obvious
global attractor: among constant-$|\B|$ states with a given mean, the
smoothest are the 1D circularly polarised waves of lowest wavenumber, so the
unconstrained descent degenerates toward 1D circular polarisation and answers
nothing. Included, the rows hold a genuinely 3D anchor fixed
(Lemma~\ref{lem:pin}) while the rest of the spectrum is allowed to relax, and
the limit of the descent is a candidate answer to the minimal-dressing (or
``quenching'') question: what is the least spectral tail compatible with a
prescribed 3D core? The module is marked experimental --- correctness over
speed, an un-jitted CG loop over jitted matvecs --- because the retraction and
the border interact in ways we have not yet characterised.
\end{remark}

\section{The continuation driver}\label{sec:driver}

The driver (\texttt{grow.py}) walks the seed amplitude upward, re-converging
after every push, and refines the grid when --- and only when --- the honest
diagnostics say the state has outgrown it. Algorithm~\ref{alg:grow} is the
loop.

\begin{algorithm}[htbp]
\caption{Carrier-free continuation (\texttt{grow.py})}
\label{alg:grow}
\small
\begin{algorithmic}[1]
\State $\B\gets\bar\B$ (uniform); refuse if \texttt{fix-mean} and
$|\bar\B|\ge1-10^{-9}$ \Comment{Corollary~\ref{cor:mean}}
\State build seed $\bb=\Trr\nabla\times\mathbf a$, $\max|\bb|=1$; record its
spec in the state metadata
\State if \texttt{freeze-top} $m>0$: freeze the $m$ dominant bins of $\hat\bb$
(plus $\kk=0$ if \texttt{fix-mean}); report the triples and the
non-coplanarity verdict \Comment{Lemma~\ref{lem:pin}}
\State $\B\gets\Trr\Pl\B$ \Comment{project once, \eqref{eq:project}}
\State $\epsilon\gets0$; $\delta\gets\delta_0$; $c\gets0$
\While{$\epsilon<\epsilon_{\max}$ and time remains}
  \State $\B_{\text{trial}}\gets\B+\delta\,\bb$ \Comment{$\bb$ rebuilt on the
  current grid from metadata}
  \State $(\B_{\text{trial}},r,\texttt{cg})\gets\textsc{GaussNewton}
  (\B_{\text{trial}})$ \Comment{\eqref{eq:step-G}, a few sweeps}
  \If{$r>r_{\text{ok}}$} \Comment{no salvage: quadratic Newton means a bad
  step is genuinely bad}
    \State $\delta\gets\delta/2$; $c\gets0$
    \If{$\delta<\delta_{\min}$} \State \textbf{stop:} \textsc{Fold candidate}
    \EndIf
    \State \textbf{continue}
  \EndIf
  \State $\B\gets\B_{\text{trial}}$; $\epsilon\gets\epsilon+\delta$;
  $c\gets c+1$
  \State record diagnostics (Section~\ref{sec:bookkeeping}) to CSV; snapshot
  every $\Delta\epsilon_{\text{snap}}$
  \If{$\texttt{gal\_tail\_rms}>\tau_{\text{tail}}$ \textbf{or}
      $\texttt{edge}>\tau_{\text{edge}}$} \Comment{honest-tail / band-edge
      triggers}
    \If{grid $=$ \texttt{grid-max}} \State \textbf{stop:}
      \textsc{Sharpening} \EndIf
    \State grid $\gets\min(1.5\times\text{grid},\texttt{grid-max})$;
      $\B\gets\textsc{ZeroPad}(\B)$; re-polish with extra sweeps and report if
      $r_{\text{ok}}$ is not met
    \State $c\gets0$
  \EndIf
  \If{$c\ge2$} \State $\delta\gets\min(1.3\,\delta,\delta_{\max})$;
  $c\gets0$ \EndIf
\EndWhile
\end{algorithmic}
\end{algorithm}

Three features of Algorithm~\ref{alg:grow} are not tuning choices.

\paragraph{The ladder ascends, never descends.} Refinement is always
zero-padding a converged state to a finer grid and re-polishing. The reverse
operation --- truncating a fine state and re-converging --- is forbidden,
because it lets the solver relax onto a different, smoother member of the
solution family: it was observed to do exactly that at $\epsilon=0.98$ on a
$32^2\times64$ ladder rung. An ascending ladder tracks one continuum object;
a descending one silently changes the subject.

\paragraph{There is no salvage branch.} With quadratic convergence
(Theorem~\ref{thm:quad}), a step that has not reached $r_{\text{ok}}$ after
the allotted sweeps is not a step that needs more sweeps; it is a step that
was too large. The response is to reject, halve $\delta$, and retry. Repeated
halving to $\delta_{\min}$ is reported as a \textsc{Fold candidate}: the
continuation cannot proceed at any resolvable step size, which is the
numerical signature of a turning point in the path.

\paragraph{Refinement triggers are mode-specific.} In Galerkin mode the
grid's top modes are identically zero, so the classical ``field content near
the Nyquist frequency'' trigger is blind. The v2 triggers are the honest tail
(\texttt{gal\_tail\_rms}, Proposition~\ref{prop:tail}) and the field's content
at the retained-band edge. Exhausting the top grid while the tail still grows
is the \textsc{Sharpening} stop: the state is demanding resolution faster than
the ladder can supply it, which is evidence about the solution, not about the
solver.

\subsection{Amplitude bookkeeping}\label{sec:bookkeeping}

The continuation parameter $\epsilon=\sum\delta$ is a \emph{path label}: the
accumulated injected seed amplitude. It is not a property of the state, and
two paths at the same $\epsilon$ are not comparable. The state properties
recorded per accepted step are the fluctuation level
$\texttt{fluct\_rms}=\mathrm{rms}\,|\B-\bar\B|$, the drift $1-|\bar\B|$, the
maximum deflection from $\bar\B$, the reversed volume fraction, the maximum
gradient, the residual $\max|\Tr q|$, the tail norms, the band-edge content,
the CG count, and the inertia ratios $\lambda_2/\lambda_1$,
$\lambda_3/\lambda_1$ of \eqref{eq:inertia}. Note that at convergence the
$\kk=0$ component of $\Tr q$ vanishes, so $\langle|\B|^2\rangle=1$ exactly and
\eqref{eq:energy} holds: $\texttt{fluct\_rms}^2=1-|\bar\B|^2$. The two
amplitude measures are therefore \emph{not} independent, and with a fixed mean
field neither of them moves at all --- which is precisely why $\epsilon$,
the gradient measures, and the tail measures are all logged separately.

\section{Practical notes}\label{sec:practical}

\paragraph{Float64 is mandatory.} The CG stop is
$\|r\|^2<10^{-26}\|r_0\|^2$, i.e.\ relative residual $10^{-13}$, and
Theorem~\ref{thm:quad} takes the Newton residual to $10^{-15}$ in four sweeps;
neither is meaningful in single precision. \texttt{jax.config} enables x64
once, in the package \texttt{\_\_init\_\_}, before anything else touches JAX.
\texttt{jax-metal} has no x64 support, so Mac GPUs are not a target; the P100's
$1{:}2$ fp64 throughput is adequate.

\paragraph{One compiled program per solve.} The sweep loop and the CG loop are
both \texttt{lax.while\_loop}s inside a single \texttt{jax.jit}, so a
multi-sweep solve runs on-device with no per-iteration host dispatch.
Code-path flags are static Python booleans (a handful of compiled variants per
grid shape); numeric knobs (\texttt{sweeps}, \texttt{cgit}, \texttt{tol}) are
traced, so changing them does not recompile. Buffer donation is not used
anywhere in v2: it deletes caller arrays, and the continuation driver keeps
the previous state alive for rejection.

\paragraph{The Nyquist convention.} Derivative wavenumbers zero the
$|k_i|=N_i/2$ bins, because $i\kk$ makes the self-conjugate Nyquist planes
anti-Hermitian and a real inverse transform discards their contribution; the
$k^2$-only quantities (the weight $\Wm^{\pm2}$, the preconditioner
\eqref{eq:pc}) keep the true Nyquist values. Mixing the two conventions
corrupts divergence residuals of states with near-Nyquist content by orders of
magnitude. In Galerkin mode the retained band excludes the Nyquist planes
anyway, but the operators are shared with the legacy collocation path.

\paragraph{Memory.} A solve holds the state, the seed and the CG workspace
(the scalar fields $\mu,r,z,p$ plus transform scratch); the dominant term is
the batched complex half-spectrum of a three-component field,
$3N_xN_y(N_z/2+1)$ complex doubles. The legacy coupled solver peaks near
$5$\,GB at $192^2\times384$, which fits a $16$\,GB P100; the $\mu$-only solver
carries one scalar unknown instead of four fields and is strictly smaller.

\paragraph{Seeds are grid-independent.} Random seed coefficients are drawn
from a generator keyed by the \emph{integer wavevector} and the run key, not
by array index, so the same key produces the same continuum field on every
grid. This is what makes a ladder rung's re-seeded push identical to the
previous rung's, up to resolution: a seed built on $24^2\times48$ is the
zero-padding of the same seed built on $16^2\times32$ wherever both resolve
it. Without this the ladder would perturb the state every time it refined.

\paragraph{Mode compatibility.} Galerkin and collocation runs are not mutually
resumable: the residuals mean different things and the CSV schemas differ.
Start fresh state and CSV files when switching.

\section*{Summary of the equations the code implements}

The hot path applies, in order: $\Pl_\Wm\vv:=\Wm^{-2}\Zf\Trr\Pl\vv$ (one
batched transform pair); the normal operator
$\Aop\mu=\Tr[\B\cdot\Pl_\Wm(\mu\B)]$, equation~\eqref{eq:step-G}; the
preconditioner $\Minv=\Tr(1+k^2)^{s}$, equation~\eqref{eq:pc}; the CG
right-hand side $-\Tr q$ for a solve step and $\Tr q+1$ for a descent step,
equation~\eqref{eq:sqp-mu}; and the sweep
$\B\leftarrow\Trr(\B+\Pl_\Wm(\mu\B))$, whose truncation is a no-op by
Theorem~\ref{thm:quad} and is kept only as a guard.

\begin{question}
The measurements of Section~\ref{sec:legacy} identify the coupled system, not
operator projection, as the source of the legacy limit cycle, but do so by
comparing two implementations that differ in more than one respect. A minimal
reproducer --- the coupled system with an unbounded CG budget, and the
$\mu$-only system with a truncated one --- would settle it.
\end{question}

\end{document}
